% !PRINTED 109
\weberpage{109}

Nun ist $a + b\varrho$ eine ganze oder gebrochene Zahl des Körpers $R(\varrho)$.
Wenn wir Zähler und Nenner dieser Zahl in ihre Primfactoren zerlegen, gemeinsame
Factoren wegheben, und wenn wir unter $q_1, q_2, \ldots$ reelle Primzahlen der
Form $3N+2$, unter $\pi_1, \pi_1', \pi_2, \pi_2', \ldots$ conjugirte Paare
imaginärer Primzahlen der Form (6) verstehen, so erhalten wir

\weq{8}{\begin{aligned}
a + b\varrho &= (-\varrho)^\lambda (\sqrt{-3})^n q_1^{s_1} q_2^{s_2} \ldots \pi_1^{t_1} \pi_1'^{t_1'} \pi_2^{t_2} \pi_2'^{t_2'} \ldots \\
a + b\varrho^2 &= (-\varrho^2)^\lambda (-\sqrt{-3})^n q_1^{s_1} q_2^{s_2} \ldots \pi_1^{t_1'} \pi_1'^{t_1} \pi_2^{t_2'} \pi_2'^{t_2} \ldots
\end{aligned}}

worin

\[
\lambda, \ n, \ s_1, \ s_2 \ldots t_1, \ t_1', \ t_2, \ t_2' \ldots
\]

ganze positive oder negative Zahlen sind.

Nun ist aber nach (3) das Product der beiden Zahlen (8) der Cubus einer
rationalen Zahl, und dies giebt, da eine rationale Zahl nur auf eine Art in
Primfactoren zerlegbar ist, die folgenden Bedingungen:

\weq{8}{\begin{aligned}
n &\equiv 0, \quad s_1 \equiv 0, \quad s_2 \equiv 0 \ldots \\
t_1+t_1' &\equiv 0, \quad t_2+t_2' \equiv 0 \ldots \quad \text{(mod 3),}
\end{aligned}} % [sic: Weber repeats equation number (8); (9) follows.]

und wir setzen also

\weq{9}{\begin{aligned}
n &= 3\mu, \quad s_1 = 3\sigma_1, \quad s_2 = 3\sigma_2 \ldots \\
t_1+t_1' &= 3\tau_1, \quad t_2+t_2' = 3\tau_2 \ldots,
\end{aligned}}

so dass $n, \sigma_1, \sigma_2 \ldots \tau_1, \tau_2 \ldots$ ganze Zahlen sind. % [sic: Weber prints n, although n=3\mu has just been set and \mu is evidently intended.]
Nun ist nach den Formeln (7):

\weq{10}{\begin{aligned}
\pi^t \pi'^{t'} &= p^{-(t+t')} (\varrho,\eta)^{3t} (\varrho^2,\eta)^{3t'} \\
\pi^{t'} \pi'^{t} &= p^{-(t+t')} (\varrho,\eta)^{3t'} (\varrho^2,\eta)^{3t}.
\end{aligned}}

Bezeichnen wir endlich noch mit $\varepsilon$ eine $9^{\text{te}}$
Einheitswurzel, so können wir

\[
-\varrho = (-\varepsilon)^3
\]

setzen, und stellen dadurch $a+b\varrho$ und $a+b\varrho^2$ als Cuben dar. Wir
erhalten also nach (2):

\weq{11}{(\varrho, x_0)^3 = (-\varepsilon)^{3\lambda} (\sqrt{-3})^{3\mu} q_1^{3\sigma_1} q_2^{3\sigma_2} \ldots p_1^{-3\tau_1} p_2^{-3\tau_2} \ldots H_1^3,}

worin

\weq{12}{H_1 = (\varrho,\eta_1)^{t_1} (\varrho^2,\eta_1)^{t_1'} (\varrho,\eta_2)^{t_2} (\varrho^2,\eta_2)^{t_2'} \ldots}

eine Kreistheilungszahl ist. Bezeichnen wir mit $H_2$ die aus $H_1$ durch
Vertauschung von $\varrho$ mit $\varrho^2$ hervorgehende Zahl, so ist $H_1 H_2$
eine rationale Zahl.

Wenn wir aus (11) die dritte Wurzel ziehen, so folgt, da hierbei noch eine
dritte Einheitswurzel $\varrho_1$ als Factor auftreten kann,

\weq{13}{(\varrho, x_0) = (-\varepsilon)^\lambda \varrho_1 (\sqrt{-3})^\mu q_1^{\sigma_1} q_2^{\sigma_2} \ldots p_1^{-\tau_1} p_2^{-\tau_2} \ldots H_1,}
%CONT>>
